\section{El cerebro del robot: el sistema dual de VLM y VLA}

El ciclo cerrado de datos resuelve «de qué se alimenta el modelo»; el cerebro del robot resuelve «cómo actúa el modelo». Aquí es también donde este episodio pasa del análisis de la industria a la estructura técnica.

Cuando el presentador pregunta «cómo se construye el cerebro», Gao Jiyang (高继扬) lo divide en dos modelos fundacionales: uno es el VLM de nivel superior, que se encarga de descomponer instrucciones, razonar y planificar tareas; el otro es el modelo fundacional de acción, el VLA, que a partir de vision y language produce action y hace que el cuerpo del robot ejecute la tarea. Esta es la estructura técnica más importante del episodio.

\begin{figure}[H]
\centering
\includegraphics[width=0.96\textwidth]{robot-brain-dual-system.png}
\caption{Estructura de sistema dual del cerebro del robot: el VLM descompone las instrucciones y el VLA produce las acciones.}
\end{figure}

\begin{knowledgebox}{Lectura de la figura: por qué separar el VLM y el VLA}
El VLM es adecuado para interpretar instrucciones ambiguas, descomponerlas lógicamente y planificar tareas de carácter más general; el VLA debe producir acciones con baja latencia y, en muchos casos, tiene que ejecutarse en el dispositivo. En los escenarios industriales, el conjunto de tareas es limitado y puede invocarse directamente la interfaz de lenguaje del VLA; en escenarios generales, como el hogar, se necesita más que el VLM participe en la descomposición de tareas complejas.
\end{knowledgebox}

\subsection{Latencia en el dispositivo e inferencia en la nube}

Gao Jiyang señala en particular que, si el modelo que ejecuta las acciones se coloca en la nube, la latencia se vuelve un problema muy difícil de resolver. Por eso el VLA, como modelo de acción, suele tener que desplegarse en el dispositivo. Que el VLM deba participar o no en tiempo real depende de la complejidad del escenario. En escenarios industriales y comerciales con veinte o treinta acciones fijas, quizá no sea necesario invocar el VLM en cada paso; en escenarios domésticos más generales, el VLM sí es un componente indispensable.

\begin{warningbox}{No hay que imaginar el «cerebro del robot» como un único modelo grande}
El cerebro del robot no es un LLM monolítico. Incluye, como mínimo, descomposición de tareas, generación de acciones, control en el dispositivo, entrada de percepción, retroalimentación de vuelta al sistema y restricciones de seguridad. Reducirlo a «conectar un modelo grande de lenguaje» subestima los problemas de despliegue y de latencia.
\end{warningbox}

\subsection{Ventajas de las empresas emergentes frente a las grandes empresas}

La sección anterior separó el VLM y el VLA; esta responde a la pregunta de la competencia: si las grandes empresas tienen modelos fundacionales, capacidad de cómputo y talento, ¿con qué puede una empresa emergente construir el cerebro del robot? La respuesta de Gao Jiyang no es afirmar de forma general que son «más rápidas», sino volver al punto de entrada de los datos.

Gao Jiyang divide los factores de éxito para construir un buen VLA en datos, algoritmos, capacidad de cómputo/infraestructura y talento. Las grandes empresas son fuertes en capacidad de cómputo, infraestructura y talento, pero con frecuencia carecen de datos reales de manipulación; las empresas emergentes que tienen capacidad de fabricar el robot completo pueden ser más fuertes en el punto de entrada de los datos. Sobre todo en China, la cadena de suministro y la capacidad de iterar el hardware facilitan que las empresas emergentes integren el robot completo, los escenarios y la recolección de datos.

\subsection{Resumen de la sección}

Lo decisivo en el cerebro del robot no es solo «quién tiene el modelo grande más potente», sino si se logra contar con un VLA ejecutable en el dispositivo, una descomposición adecuada a cargo del VLM en el nivel superior, un ciclo cerrado de datos reales y la capacidad de desplegar el robot completo.
